\rtmtight
\begin{tblr}{width=\textwidth,colspec={|Q[c,m,2.1cm]|X[l,m]|},hlines,vlines,hspan=minimal,rowsep=1.5pt,colsep=3pt,row{1}={bg=headgray,font=\bfseries}}
\SetCell{c,m}\hfil\logo\hfil & \textbf{POLITEKNIK NEGERI MALANG}\par\textbf{JURUSAN TEKNOLOGI INFORMASI}\par\textbf{PROGRAM STUDI: D4 TEKNIK INFORMATIKA} \\
\SetCell[c=2]{l} \textbf{RENCANA TUGAS MAHASISWA (RTM) DAN RUBRIK PENILAIAN} & \\
Nama Mata Kuliah & Bahasa Inggris 2 \\
Kode Mata Kuliah & RTI253009 \quad \textbf{sks / Jam perminggu:} 2 SKS/4 jam \quad \textbf{Semester:} 3 \\
Dosen Pengampu & Tim Pengajar Program Studi D4 TI \\
\end{tblr}
\assessmentblock{Tugas Writing dan Reading}{Case Method dan praktik komunikasi}{SCPMK0102-02701, SCPMK0303-02702}{Mahasiswa menyelesaikan tugas membaca teks teknis IT, menulis laporan akademik dan teknis, serta berlatih keterampilan komunikasi lisan sesuai topik mingguan.}{Membaca teks yang ditugaskan, mengerjakan latihan menulis atau berbicara, mendokumentasikan hasil, dan mengumpulkan evidence sesuai petunjuk.}{Tulisan akademik/teknis, ringkasan, laporan, rekaman presentasi/diskusi, dan/atau CV dan email profesional.}{Kualitas tulisan akademik dan teknis & 12 & Laporan, esai, atau teks teknis ditulis dengan struktur yang benar, kosakata teknis yang tepat, dan tata bahasa yang baik. & Tulisan memiliki struktur yang cukup baik tetapi kosakata atau tata bahasa masih perlu perbaikan. & Tulisan tidak terstruktur, kosakata terbatas, atau tata bahasa banyak kesalahan. \\ \\
Kemampuan membaca dan merangkum teks teknis & 8 & Teks teknis IT dipahami dengan baik, gagasan utama diidentifikasi tepat, dan ringkasan disusun secara koheren. & Sebagian teks dipahami dan dirangkum tetapi detail atau koherensi masih kurang. & Teks tidak dipahami atau ringkasan tidak mencerminkan isi. \\ \\
Kualitas komunikasi lisan dan presentasi & 10 & Presentasi atau diskusi disampaikan dengan lancar, kosakata teknis digunakan tepat, dan ide dijelaskan secara terstruktur. & Komunikasi cukup dipahami tetapi ada gangguan kelancaran atau penggunaan kosakata yang kurang tepat. & Komunikasi tidak jelas atau ide tidak tersampaikan dengan baik. \\}

\assessmentblock{Kuis Reading, Vocabulary, Listening, dan Speaking}{Kuis}{SCPMK0102-02701, SCPMK0303-02702}{Kuis tertulis dan lisan untuk mengukur penguasaan kosakata teknis IT (Kuis 1) dan kemampuan mendengarkan serta berbicara dalam bahasa Inggris (Kuis 2).}{Menjawab soal pilihan ganda dan/atau uraian terkait bacaan teknis, kosakata, mendengarkan presentasi, dan respons lisan spontan.}{Jawaban tertulis kuis dan/atau rekaman respons lisan.}{Penguasaan kosakata teknis IT & 3 & Kosakata teknis IT digunakan dan dipahami dengan tepat dalam konteks kalimat dan paragraf teknis. & Sebagian kosakata dipahami tetapi penggunaannya dalam konteks masih terbatas. & Kosakata teknis IT tidak dikuasai atau banyak kesalahan. \\ \\
Pemahaman bacaan dan mendengarkan & 3.5 & Teks bacaan dan audio teknis dipahami dengan akurat, pertanyaan dijawab dengan benar dan lengkap. & Sebagian besar isi dipahami tetapi beberapa detail atau inferensi masih keliru. & Pemahaman bacaan atau mendengarkan sangat terbatas. \\ \\
Kemampuan berbicara spontan & 3.5 & Respons lisan diberikan dengan tepat, lancar, dan menggunakan bahasa Inggris yang grammatically acceptable. & Respons lisan dapat dipahami tetapi ada kekeliruan grammar atau kelancaran yang mengganggu. & Respons lisan tidak dapat dipahami atau sangat terbatas. \\}

\assessmentblock{UTS Reading Comprehension dan Writing}{UTS berbasis tes tertulis}{SCPMK0102-02701, SCPMK0303-02702}{Evaluasi tengah semester untuk mengukur kemampuan membaca teks teknis IT, menulis secara akademik, dan berkomunikasi lisan dalam bahasa Inggris.}{Membaca teks teknis yang disediakan, menjawab pertanyaan pemahaman, menulis paragraf atau esai, dan memberikan respons lisan singkat.}{Lembar jawaban tertulis dan rekaman respons lisan.}{Pemahaman membaca teks teknis & 6 & Pertanyaan tentang teks teknis dijawab dengan akurat, gagasan utama, detail, dan inferensi diidentifikasi dengan benar. & Sebagian besar pertanyaan dijawab benar tetapi beberapa inferensi atau detail masih keliru. & Pemahaman teks teknis sangat terbatas atau banyak jawaban keliru. \\ \\
Kualitas penulisan & 5 & Paragraf atau esai ditulis dengan struktur yang benar, kosakata tepat, argumen koheren, dan tata bahasa yang baik. & Tulisan cukup dapat dipahami tetapi struktur atau tata bahasa masih perlu perbaikan yang signifikan. & Tulisan tidak terstruktur atau banyak kesalahan tata bahasa yang mengganggu pemahaman. \\ \\
Respons komunikatif lisan & 4 & Respons lisan diberikan dengan tepat, terstruktur, dan menggunakan bahasa Inggris yang jelas dan dapat dipahami. & Respons dapat dipahami tetapi kelancaran atau struktur masih terbatas. & Tidak dapat berkomunikasi secara efektif dalam bahasa Inggris lisan. \\}

\assessmentblock{PBL English Technical Presentation Project}{Project Based Learning}{SCPMK0303-02702}{Mahasiswa menyiapkan dan menyampaikan presentasi teknis lengkap dalam bahasa Inggris tentang topik teknologi informasi, dilengkapi slide dan sesi tanya jawab.}{Memilih topik TI, melakukan riset, merancang slide, menyiapkan skrip, mempraktikkan delivery, dan mempresentasikan di depan kelas.}{Slide presentasi, rekaman presentasi, dan/atau laporan singkat topik.}{Kualitas konten teknis presentasi & 9 & Topik teknis TI disampaikan secara akurat, mendalam, dan menggunakan kosakata teknis yang tepat serta relevan. & Konten teknis cukup akurat tetapi kurang mendalam atau kosakata teknis masih terbatas. & Konten teknis tidak akurat, tidak relevan dengan bidang TI, atau sangat dangkal. \\ \\
Kelancaran dan delivery presentasi & 8 & Presentasi disampaikan dengan lancar, percaya diri, eye contact baik, intonasi yang tepat, dan waktu terkelola. & Presentasi dapat dipahami tetapi kelancaran, percaya diri, atau pengelolaan waktu masih perlu peningkatan. & Presentasi tidak lancar, sulit dipahami, atau tidak dapat diselesaikan. \\ \\
Kualitas materi dan desain slide & 5 & Slide dirancang dengan baik, visual mendukung konten, alur presentasi logis, dan mudah diikuti oleh audiens. & Slide cukup baik tetapi terlalu padat teks atau desain kurang mendukung pemahaman audiens. & Slide tidak informatif, tidak relevan, atau tidak mendukung penyampaian konten teknis. \\}

\assessmentblock{UAS Final Presentation dan Writing Test}{UAS berbasis tes dan presentasi}{SCPMK0102-02701, SCPMK0303-02702}{Ujian akhir semester yang mengukur kemampuan menulis teks teknis secara akademik dan mempresentasikan topik TI dalam bahasa Inggris.}{Mengerjakan tes menulis tentang topik teknis yang diberikan dan menyampaikan presentasi final dalam bahasa Inggris.}{Tulisan akhir teknis dan rekaman presentasi final.}{Kualitas tulisan akhir & 10 & Teks teknis ditulis dengan struktur yang baik, kosakata teknis tepat, argumen koheren, dan tata bahasa yang benar. & Tulisan cukup baik tetapi ada kekurangan pada struktur argumen atau ketepatan tata bahasa. & Tulisan tidak terstruktur atau banyak kesalahan yang mengganggu pemahaman. \\ \\
Presentasi final teknis & 8 & Topik teknis dipresentasikan secara menyeluruh, akurat, dan dengan delivery yang meyakinkan dalam bahasa Inggris. & Presentasi cukup baik tetapi ada kekurangan pada kedalaman konten atau delivery. & Presentasi tidak memadai, konten tidak akurat, atau tidak dapat diselesaikan. \\ \\
Kemampuan Q\&A dalam bahasa Inggris & 5 & Pertanyaan dipahami dan dijawab dengan tepat menggunakan bahasa Inggris yang baik, lancar, dan relevan. & Pertanyaan dijawab tetapi respons masih terbatas atau menggunakan bahasa yang kurang tepat. & Tidak dapat menjawab pertanyaan dalam bahasa Inggris atau jawaban tidak relevan. \\}

\begin{tblr}{width=\textwidth,colspec={|Q[l,m,3.2cm]|X[l,m]|},hlines,vlines,hspan=minimal,row{1-3}={bg=headgray},rowsep=1.5pt,colsep=3pt}
Jadwal Pelaksanaan: & Minggu 1--17 sesuai RPS. Setiap evaluasi memiliki RTM dan rubrik multi-kriteria. \\
Lain-Lain Yang Diperlukan: & LMS, sumber bacaan teknis IT berbahasa Inggris, perangkat audio/video, dan perangkat presentasi. \\
Pustaka: & Mengacu pustaka utama dan pendukung pada bagian RPS. \\
\end{tblr}
